%% RL for LLMs: Pen-and-Paper Exercises
%% Papers listed from: https://algoroxyolo.github.io/blog/2026/rl-reading-list/
%% 96 papers · 5 categories · 24 sub-topics
%% TODO markers indicate exercises that still need to be written.

\documentclass[11pt,a4paper]{book}

\usepackage[T1]{fontenc}
\usepackage[utf8]{inputenc}
\usepackage{amsmath,amssymb,amsthm}
\usepackage{enumitem}
\usepackage{hyperref}
\usepackage{geometry}
\usepackage{xcolor}
\usepackage{booktabs}
\geometry{margin=2.5cm}
\hypersetup{colorlinks=true, linkcolor=blue, urlcolor=teal}

% ── priority labels ────────────────────────────────────────────────────────────
\newcommand{\pA}{\textbf{[A]}}     % deep read
\newcommand{\pAm}{\textbf{[A-]}}   % method + experiments
\newcommand{\pBp}{\textbf{[B+]}}   % problem + conclusion
\newcommand{\pB}{\textbf{[B]}}     % skim
\newcommand{\pC}{\textbf{[C]}}     % core idea only
\newcommand{\NV}{\textcolor{teal}{\textsc{nv}}}  % NVIDIA paper

% ── exercise environment ───────────────────────────────────────────────────────
\newcounter{exno}[chapter]
\renewcommand{\theexno}{\thechapter.\arabic{exno}}

\newenvironment{exercise}[1][]{%
  \refstepcounter{exno}%
  \par\smallskip\noindent
  \textbf{Exercise~\theexno%
  \ifx&#1&\else\ — #1\fi.}\quad
}{%
  \par\smallskip
}

% ── TODO macro ─────────────────────────────────────────────────────────────────
\newcommand{\TODO}[1][\relax]{%
  \colorbox{yellow!40}{%
    \textbf{TODO}%
    \ifx#1\relax\else: #1\fi
  }%
}

% ── paper entry macro ─────────────────────────────────────────────────────────
% \paper{year}{title}{notes}{priority}
\newcommand{\paper}[4]{%
  \item[\textbullet] \textit{(#1)} \textbf{#2}%
  \ifx&#3&\else\ — {\small #3}\fi
  \quad #4%
}

% ── section note macro ────────────────────────────────────────────────────────
\newcommand{\sectionnote}[1]{\par{\small\textit{#1}}\par\medskip}

% ── title page ────────────────────────────────────────────────────────────────
\title{\Huge\bfseries RL for LLMs\\[0.5ex]
       \Large Pen-and-Paper Exercises\\[1ex]
       \normalsize Based on the 96-paper reading list\\
       \texttt{algoroxyolo.github.io/blog/2026/rl-reading-list/}}
\author{}
\date{March 2026}

\begin{document}
\maketitle
\tableofcontents

% ══════════════════════════════════════════════════════════════════════════════
% HOW TO USE
% ══════════════════════════════════════════════════════════════════════════════
\chapter*{How to Use This Book}
Each chapter corresponds to one of the five categories in the reading list.
Papers are listed with their year, title, brief notes, and priority label.
Below each paper (or small cluster of papers) there is one or more \textbf{Exercise}
slot marked \TODO{} — that is where an exercise question needs to be written.

Priority labels: \pA\ deep read · \pAm\ method + experiments · \pBp\ problem + conclusion
· \pB\ skim · \pC\ core idea only · \NV\ NVIDIA paper.

% ══════════════════════════════════════════════════════════════════════════════
\chapter{Algorithms}
\sectionnote{26 papers across four sub-topics.}
% ──────────────────────────────────────────────────────────────────────────────
\section{Core RL: PPO $\to$ GRPO $\to$ DAPO $\to$ CISPO $\to$ MaxRL $\to$ DPPO}
\sectionnote{13 papers.  Start with PPO, end with DPPO; each paper refines one design choice.}

\begin{description}
\paper{2017}{PPO}{Know ratio clipping + trust region}{\pB}
\paper{2022}{InstructGPT}{RM + PPO + KL penalty; shared with \S3.1}{\pAm}
\paper{2024}{RLOO / Revisiting REINFORCE}{}{\pAm}
\paper{2024}{DeepSeekMath (GRPO origin)}{Start your interview story here}{\pA}
\paper{2025}{DeepSeek-R1}{}{\pA}
\paper{2025}{DAPO}{Recipe + system-level scaling}{\pA}
\paper{2025}{GSPO}{}{\pAm}
\paper{2025}{SAPO}{}{\pB}
\paper{2025}{MiniMax-M1 / CISPO}{Changed clipping target — key interview talking point; also notable for long-context RL}{\pA}
\paper{2025}{PPO Collapse in Long-CoT}{}{\pBp}
\paper{2025}{Dr.\ GRPO}{Fixes length bias + std normalisation}{\pBp}
\paper{2026}{MaxRL}{pass@k objective; RL--MLE continuum}{\pAm}
\paper{2026}{DPPO}{TV/KL divergence trust region}{\pBp}
\end{description}

\begin{exercise}[PPO clipping]
Let $r_t(\theta)=\dfrac{\pi_\theta(a_t\mid s_t)}{\pi_{\theta_{\text{old}}}(a_t\mid s_t)}$ be the probability ratio and $\hat{A}_t$ the estimated advantage.
\begin{enumerate}[label=(\alph*)]
  \item Write the PPO clipped surrogate objective
        \[L^{\text{CLIP}}(\theta)=\mathbb{E}_t\!\left[\min\!\left(r_t(\theta)\hat{A}_t,\;
        \operatorname{clip}(r_t(\theta),1{-}\varepsilon,1{+}\varepsilon)\hat{A}_t\right)\right].\]
        Identify the role of each argument of $\min$.
  \item Prove that when $\hat{A}_t>0$ and $r_t(\theta)>1+\varepsilon$,
        the gradient $\nabla_\theta L^{\text{CLIP}}(\theta)=\mathbf{0}$ (with respect to the clipped term).
        (\emph{Hint}: evaluate $\operatorname{clip}(r_t,1-\varepsilon,1+\varepsilon)$ in this region.)
  \item By an analogous argument, show that when $\hat{A}_t<0$ and $r_t(\theta)<1-\varepsilon$
        the clipped term's gradient also vanishes.
  \item For $\varepsilon=0.2$ and a fixed $\hat{A}_t=1$, sketch $L^{\text{CLIP}}$ as a function of
        $r_t(\theta)$ over $[0.5,\,2.0]$.  Mark the two flat regions and the linear segment.
  \item Explain intuitively why these flat regions implement a soft trust-region constraint
        without requiring a second-order optimisation step.
\end{enumerate}
\end{exercise}

\begin{exercise}[PPO: TRPO connection and KL penalty]
The PPO paper (Schulman et al., 2017) presents the clipped objective as a simpler alternative
to TRPO.  It also proposes a KL-penalty fallback.
\begin{enumerate}[label=(\alph*)]
  \item Write the TRPO constrained optimisation problem; identify the step that requires
        computing (or approximating) the Fisher information matrix, and state its
        memory complexity for a model with $n$ parameters.
  \item Write the PPO KL-penalty objective
        $L^{\text{KL}}(\theta)=\mathbb{E}_t\!\left[r_t(\theta)\hat{A}_t\right]
        -\beta\,\overline{\mathrm{KL}}\!\left[\pi_{\theta_{\text{old}}},\pi_\theta\right]$
        (where $\overline{\mathrm{KL}}$ denotes the mean KL divergence over the batch).
        State which KL direction is used and why.
  \item Describe the \emph{adaptive} KL coefficient scheme: given a target
        $d_{\text{targ}}$, what adjustment is made to $\beta$ when the observed KL is
        (i) greater than $1.5\,d_{\text{targ}}$, and (ii) less than $d_{\text{targ}}/1.5$?
  \item Give one advantage and one disadvantage of the KL-penalty approach compared
        to the clipped approach for practical large-scale training.
\end{enumerate}
\end{exercise}

\begin{exercise}[Generalised Advantage Estimation in PPO]
PPO estimates advantages using GAE$(\gamma,\lambda)$ (Schulman et al., 2016).
\begin{enumerate}[label=(\alph*)]
  \item Define the one-step TD residual
        $\delta_t^V = r_t + \gamma V(s_{t+1}) - V(s_t)$
        and write the GAE estimator as an infinite weighted sum:
        $\hat{A}_t^{\text{GAE}(\gamma,\lambda)}
         =\sum_{k=0}^{\infty}(\gamma\lambda)^k \delta_{t+k}^V$.
  \item Show that $\lambda=0$ recovers the one-step TD advantage
        $r_t+\gamma V(s_{t+1})-V(s_t)$,
        and that $\lambda=1$ recovers
        $\sum_{k=0}^{\infty}\gamma^k r_{t+k} - V(s_t)$
        (Monte-Carlo return minus baseline).
  \item Explain the bias-variance trade-off: which endpoint ($\lambda=0$ or $\lambda=1$)
        has higher variance?  Which has higher bias?  Justify both answers.
  \item For a two-step trajectory with $\gamma=0.99$, $\lambda=0.95$,
        $\delta_0^V=1.0$, $\delta_1^V=-0.5$, compute $\hat{A}_0^{\text{GAE}}$
        to three significant figures.
        (\emph{Note}: truncate at $k=1$ since the episode ends after step 1.)
\end{enumerate}
\end{exercise}

\begin{exercise}[PPO actor-critic algorithm]
PPO trains a shared actor-critic network using $K$ epochs of minibatch gradient ascent
on each batch of $T$ timesteps collected per actor.
\begin{enumerate}[label=(\alph*)]
  \item Write the full per-step loss
        $L^{\text{PPO}}(\theta)
         = L^{\text{CLIP}}(\theta)
           - c_1 L^{\text{VF}}(\theta)
           + c_2 S[\pi_\theta](s_t)$,
         defining $L^{\text{VF}}$ (the value-function loss) and $S[\pi_\theta]$ (the entropy bonus).
        Why is the value loss subtracted rather than added?
  \item Explain why performing $K>1$ gradient steps on the \emph{same} rollout batch
        is generally invalid under the on-policy policy-gradient theorem, and how
        the clipping mechanism makes it approximately safe.
  \item After three minibatch updates the probability ratio for a particular
        $(s_t,a_t)$ pair has grown to $r_t(\theta)=2.4$.
        With $\varepsilon=0.2$ and $\hat{A}_t=0.6$:
        \begin{enumerate}[label=(\roman*)]
          \item Compute the unclipped surrogate value $r_t(\theta)\hat{A}_t$.
          \item Compute the clipped surrogate value
                $\operatorname{clip}(r_t(\theta),1-\varepsilon,1+\varepsilon)\hat{A}_t$.
          \item Which value does $L^{\text{CLIP}}$ use?  Justify.
        \end{enumerate}
  \item State the early-stopping heuristic the paper recommends
        to prevent the ratio from drifting too far within the $K$ epochs.
\end{enumerate}
\end{exercise}

\begin{exercise}[From PPO to GRPO]
\TODO{Describe the three algorithmic changes GRPO makes relative to PPO; for each change, state the motivation from the DeepSeekMath paper.}
\end{exercise}

\begin{exercise}[DAPO entropy bonus]
\TODO{}
\end{exercise}

\begin{exercise}[CISPO clipping target]
\TODO{What quantity does CISPO clip instead of the probability ratio?  Write the modified objective.}
\end{exercise}

\begin{exercise}[Dr.\ GRPO length bias]
\TODO{}
\end{exercise}

\begin{exercise}[MaxRL objective]
\TODO{Write the MaxRL objective in terms of pass@k; show how it recovers standard REINFORCE as a special case.}
\end{exercise}

\begin{exercise}[DPPO trust region]
\TODO{Compare the DPPO TV/KL trust region to PPO's clipping; when is the TV constraint tighter?}
\end{exercise}

% ──────────────────────────────────────────────────────────────────────────────
\section{Scaling Laws \& Meta}
\sectionnote{5 papers.}

\begin{description}
\paper{2025}{Qwen3 Technical Report}{Thinking/non-thinking unified framework + thinking budget + distillation recipe}{\pA}
\paper{2025}{ScaleRL / The Art of Scaling RL Compute}{}{\pA}
\paper{2025}{Scaling Behaviors of LLM RL Post-Training}{Power-law: model scale $\times$ data $\times$ compute}{\pB}
\paper{2025}{ProRL}{Steps scaling dimension; prerequisite to ScaleRL}{\pBp\NV}
\paper{2025}{BroRL}{Rollouts scaling dimension; complementary to ProRL}{\pBp\NV}
\end{description}

\begin{exercise}[RL compute scaling]
\TODO{Sketch the power-law relationship between RL compute and benchmark score described in ScaleRL; label the three axes (model scale, data, compute).}
\end{exercise}

\begin{exercise}[ProRL vs BroRL]
\TODO{ProRL scales the \emph{steps} dimension; BroRL scales the \emph{rollouts} dimension.  Explain what each means concretely and why they are complementary.}
\end{exercise}

\begin{exercise}[Thinking budget]
\TODO{Describe the Qwen3 thinking budget mechanism; formulate it as a constrained optimisation problem.}
\end{exercise}

% ──────────────────────────────────────────────────────────────────────────────
\section{Cascade \& Stage-wise RL (Nemotron)}
\sectionnote{4 papers.}

\begin{description}
\paper{2025}{AceReason-Nemotron 1.1}{}{\pA\NV}
\paper{2025}{Nemotron-Cascade 1}{Foundational motivation for Cascade 2}{\pAm\NV}
\paper{2026}{Nemotron 3 Super Technical Report}{LatentMoE + NVFP4 + multi-env simultaneous RL + reasoning budget control}{\pAm\NV}
\paper{2026}{Nemotron-Cascade 2}{IMO/IOI/ICPC gold medals; 30B MoE; on-policy distillation}{\pA\NV}
\end{description}

\begin{exercise}[Stage-wise RL curriculum]
\TODO{What is the motivation for stage-wise RL in Nemotron-Cascade?  List the stages and the reward signal used at each stage.}
\end{exercise}

\begin{exercise}[On-policy distillation in Cascade 2]
\TODO{}
\end{exercise}

% ──────────────────────────────────────────────────────────────────────────────
\section{Distillation (on-policy / self / context)}
\sectionnote{4 papers.}

\begin{description}
\paper{2025}{On-Policy Distillation (Thinking Machines)}{Clearest explanation of compute tradeoffs}{\pAm}
\paper{2026}{MiMo-V2-Flash}{Multi-teacher on-policy distillation; core Cascade 2 technique}{\pAm}
\paper{2026}{Self-Distilled Reasoner}{Teacher = self with privileged information}{\pBp}
\paper{2026}{On-Policy Context Distillation}{}{\pB}
\end{description}

\begin{exercise}[On-policy distillation compute tradeoff]
\TODO{Using the notation from the Thinking Machines paper, compare the FLOPs of RL vs on-policy distillation at the same final accuracy target.}
\end{exercise}

\begin{exercise}[Multi-teacher distillation]
\TODO{}
\end{exercise}

% ══════════════════════════════════════════════════════════════════════════════
\chapter{Reward Modeling}
\sectionnote{20 papers across four sub-topics.}
% ──────────────────────────────────────────────────────────────────────────────
\section{Generative Reward Models}
\sectionnote{5 papers.}

\begin{description}
\paper{2024}{Generative Verifiers: RM as Next-Token Prediction}{GRM origin; yes/no token + majority vote; 346 citations}{\pA}
\paper{2025}{HelpSteer3-Preference}{40k samples; NVIDIA latest; RM-Bench SOTA}{\pA\NV}
\paper{2025}{DeepSeek-GRM / SPCT}{SPCT online RL for GRM training; 193 citations}{\pA}
\paper{2025}{RM-R1: Reward Modeling as Reasoning}{Reason first, then score; 101 citations}{\pAm}
\paper{2026}{P-GenRM (ICLR 2026 Oral)}{Personalized GRM — directly relevant to persona research}{\pAm}
\end{description}

\begin{exercise}[Generative verifier decoding]
\TODO{A generative RM produces a yes/no token.  Show how majority-vote aggregation over $n$ samples gives an unbiased estimate of the true reward.}
\end{exercise}

\begin{exercise}[SPCT training loop]
\TODO{Describe the online RL loop used to train the DeepSeek-GRM scorer (SPCT); write the reward definition.}
\end{exercise}

\begin{exercise}[RM-R1 reasoning chain]
\TODO{}
\end{exercise}

% ──────────────────────────────────────────────────────────────────────────────
\section{Process Reward Models (PRM)}
\sectionnote{7 papers.}

\begin{description}
\paper{2024}{PAV / Rewarding Progress}{ICLR; process reward = advantage; 225 citations}{\pAm}
\paper{2025}{Lessons of Developing PRMs}{Qwen/Tongyi; PRM training practice: what works, annotation, stability}{\pAm}
\paper{2025}{ThinkPRM}{Long CoT verifier; 58 citations}{\pAm}
\paper{2025}{GenPRM}{PRM test-time generative reasoning; 21 citations}{\pAm}
\paper{2025}{ReasonFlux-PRM}{NeurIPS 2025 Spotlight; trajectory-aware step + trajectory dual supervision}{\pAm}
\paper{2025}{R-PRM}{}{\pBp}
\paper{2026}{PRISM}{}{\pBp}
\end{description}

\begin{exercise}[PRM vs ORM]
\TODO{Define process reward model (PRM) and outcome reward model (ORM); give one advantage and one disadvantage of each.}
\end{exercise}

\begin{exercise}[PAV: process reward as advantage]
\TODO{Show that the PAV reward signal is equivalent to the advantage function under certain assumptions.}
\end{exercise}

\begin{exercise}[Trajectory-aware supervision in ReasonFlux-PRM]
\TODO{}
\end{exercise}

% ──────────────────────────────────────────────────────────────────────────────
\section{Rubrics-as-Rewards}
\sectionnote{4 papers.}

\begin{description}
\paper{2025}{Rubrics as Rewards (RaR)}{Rubrics structure subjective preferences; 111 citations}{\pAm}
\paper{2025}{OpenRubrics}{}{\pBp}
\paper{2026}{Rubric-ARM}{Alternating optimisation of rubric generator + judge}{\pBp}
\paper{2026}{Golden Goose}{Non-verifiable $\to$ verifiable bridge; high value for research narrative}{\pAm\NV}
\end{description}

\begin{exercise}[Rubric design]
\TODO{Design a three-criteria rubric for evaluating a creative writing response; explain how each criterion maps to a numerical reward.}
\end{exercise}

\begin{exercise}[Rubric-ARM alternating optimisation]
\TODO{Write the two optimisation steps in Rubric-ARM; identify which parameters are frozen in each step.}
\end{exercise}

\begin{exercise}[Golden Goose verifiability bridge]
\TODO{}
\end{exercise}

% ──────────────────────────────────────────────────────────────────────────────
\section{RM Evaluation \& Benchmarks}
\sectionnote{4 papers.}

\begin{description}
\paper{2024}{How to Evaluate RM for RLHF / PPE}{Offline-online gap; 61 citations}{\pAm}
\paper{2024}{HelpSteer3 dataset}{Feedback + edit; inference-time scaling}{\pBp\NV}
\paper{2025}{RewardBench 2}{Reward hacking; 63 citations}{\pBp}
\paper{2026}{Long-form RewardBench}{}{\pB}
\end{description}

\begin{exercise}[Offline-online reward model gap]
\TODO{Explain why a reward model that scores high on an offline benchmark can fail in the online RL loop; propose one mitigation.}
\end{exercise}

\begin{exercise}[Reward hacking in RewardBench 2]
\TODO{}
\end{exercise}

% ══════════════════════════════════════════════════════════════════════════════
\chapter{RLHF \& Preference Optimisation}
\sectionnote{13 papers across three sub-topics.}
% ──────────────────────────────────────────────────────────────────────────────
\section{Foundations}
\sectionnote{3 papers.}

\begin{description}
\paper{2022}{InstructGPT}{Shared with \S1.1}{\pAm}
\paper{2022}{Bai et al.\ / HH-RLHF}{Safety alignment; HH-RLHF dataset}{\pBp}
\paper{2023}{DPO}{Know why it eventually fell short}{\pAm}
\end{description}

\begin{exercise}[RLHF three-phase pipeline]
\TODO{Describe the three phases of InstructGPT training (SFT, RM, RL); state the loss function used in each phase.}
\end{exercise}

\begin{exercise}[DPO derivation]
\TODO{Derive the DPO objective from the RLHF objective by solving the constrained optimisation problem analytically.}
\end{exercise}

\begin{exercise}[DPO limitations]
\TODO{List two failure modes of offline DPO that motivate online methods.}
\end{exercise}

% ──────────────────────────────────────────────────────────────────────────────
\section{From Simplification to Online RL}
\sectionnote{4 papers.}

\begin{description}
\paper{2024}{SimPO}{DPO simplification endpoint}{\pB}
\paper{2024}{Online Iterative RLHF}{Turning point: online $\gg$ offline DPO}{\pAm}
\paper{2025}{Asynchronous RLHF (ICLR 2025)}{}{\pBp}
\paper{2025}{OLMo 3 (SFT-DPO-RL-RL Zero)}{Fully transparent \& reproducible; all data/code/checkpoints}{\pA}
\end{description}

\begin{exercise}[SimPO reference-free objective]
\TODO{Write the SimPO loss; identify the term that replaces the reference model log-probability in DPO.}
\end{exercise}

\begin{exercise}[Online vs offline preference learning]
\TODO{Sketch the data distribution shift argument that explains why online RLHF outperforms offline DPO.}
\end{exercise}

\begin{exercise}[OLMo 3 training stages]
\TODO{}
\end{exercise}

% ──────────────────────────────────────────────────────────────────────────────
\section{Multi-turn / Social / Creative RLHF}
\sectionnote{6 papers.}

\begin{description}
\paper{2025}{PPP: Proactive and Personalized Agents}{Three-objective RL; UserVille benchmark}{\pAm}
\paper{2025}{RL from User Conversations}{Persona-conditioned rewards}{\pAm}
\paper{2025}{RLMR: Mixed Rewards for Creative Writing}{Subjective aesthetic RM + objective constraint verifier}{\pB}
\paper{2026}{HER: RL for Role-playing}{Dual-layer thinking; relevant to InCharacter work}{\pAm}
\paper{2026}{OMAR (One Model All Roles)}{}{\pB}
\paper{2026}{Social-R1}{}{\pB}
\end{description}

\begin{exercise}[Multi-objective RL for personalisation]
\TODO{PPP optimises three objectives simultaneously.  Write the scalarised reward; discuss the choice of weights.}
\end{exercise}

\begin{exercise}[Persona-conditioned rewards]
\TODO{}
\end{exercise}

\begin{exercise}[Mixed rewards for creative writing]
\TODO{RLMR combines a subjective aesthetic RM with an objective constraint verifier.  Propose a concrete weighting scheme and justify it.}
\end{exercise}

% ══════════════════════════════════════════════════════════════════════════════
\chapter{Systems}
\sectionnote{17 papers across three sub-topics.}
% ──────────────────────────────────────────────────────────────────────────────
\section{Sync vs Async Training}
\sectionnote{6 papers.}

\begin{description}
\paper{2025}{Magistral}{}{\pA}
\paper{2025}{AReaL}{}{\pAm}
\paper{2025}{slime (SGLang RL blog)}{}{\pAm}
\paper{2025}{A-3PO}{Decoupled PPO; independent dimension}{\pBp}
\paper{2026}{StaleFlow}{}{\pBp}
\paper{2026}{GAC}{}{\pBp}
\end{description}

\begin{exercise}[Synchronous vs asynchronous rollout]
\TODO{Draw a timeline diagram for synchronous and asynchronous rollout collection; label the idle time in each.}
\end{exercise}

\begin{exercise}[A-3PO decoupled PPO]
\TODO{Explain the decoupling in A-3PO; identify which components run on separate hardware and why.}
\end{exercise}

\begin{exercise}[Staleness in async RL]
\TODO{}
\end{exercise}

% ──────────────────────────────────────────────────────────────────────────────
\section{Train-Inference Mismatch \& RL Stability}
\sectionnote{9 papers.}

\begin{description}
\paper{2025}{Stabilizing RL with LLMs}{IS correction + Routing Replay theory}{\pAm}
\paper{2025}{Give Me FP32 or Give Me Death}{}{\pAm}
\paper{2025}{Defeating Nondeterminism (Thinking Machines)}{}{\pAm}
\paper{2025}{SGLang Deterministic Inference}{}{\pAm}
\paper{2025}{R3: MoE Routing Replay}{Theoretically explained by Stabilizing RL paper}{\pBp}
\paper{2025}{FP16 Mismatch}{Train/infer path consistency matters more than precision per se}{\pBp}
\paper{2025}{TP Sizes Determinism}{}{\pAm}
\paper{2025}{Unified FP8 (SGLang blog)}{}{\pAm}
\paper{2025}{DeepSeek-V3.2}{}{\pB}
\end{description}

\begin{exercise}[Importance sampling correction]
\TODO{Write the importance-sampling corrected policy gradient for the case where rollouts are collected with a policy $\pi_{\text{old}}$ but updated with $\pi_\theta$.}
\end{exercise}

\begin{exercise}[Numerical precision and reward variance]
\TODO{Explain why switching from FP32 to FP16 during inference can introduce instability in the RL training signal.}
\end{exercise}

\begin{exercise}[MoE routing replay]
\TODO{What is the routing replay trick in R3?  Why does it reduce train-inference mismatch?}
\end{exercise}

\begin{exercise}[Sources of nondeterminism]
\TODO{List four sources of nondeterminism in a distributed LLM training run; for each, state a mitigation.}
\end{exercise}

% ──────────────────────────────────────────────────────────────────────────────
\section{Long-Context RL}
\sectionnote{2 papers.}

\begin{description}
\paper{2025}{Context-Folding}{Long-horizon agent context compression/folding management}{\pBp}
\paper{2025}{QwenLong-L1.5}{}{\pB}
\end{description}

\begin{exercise}[Context compression for long-horizon agents]
\TODO{Describe the context-folding mechanism; formulate the compression problem as a maximum-coverage sub-sequence selection.}
\end{exercise}

% ══════════════════════════════════════════════════════════════════════════════
\chapter{Tasks \& Agents}
\sectionnote{21 papers across six sub-topics.}
% ──────────────────────────────────────────────────────────────────────────────
\section{Coding Agents}
\sectionnote{2 papers.}

\begin{description}
\paper{2026}{Qwen3-Coder-Next}{}{\pAm}
\paper{2026}{SWE-Master}{}{\pBp}
\end{description}

\begin{exercise}[Coding agent reward design]
\TODO{Propose a reward function for a coding agent that passes unit tests; discuss how to handle partial credit.}
\end{exercise}

% ──────────────────────────────────────────────────────────────────────────────
\section{Deep Research \& Search Agents}
\sectionnote{4 papers.}

\begin{description}
\paper{2025}{Search-R1}{790 citations}{\pA}
\paper{2025}{Tongyi DeepResearch}{}{\pAm}
\paper{2026}{Yunque DeepResearch}{}{\pBp}
\paper{2026}{How to Train DR Agent}{}{\pBp}
\end{description}

\begin{exercise}[Search-R1 reward signal]
\TODO{Describe the reward signal used to train Search-R1; compare it to the reward used for pure-reasoning (no search) GRPO.}
\end{exercise}

\begin{exercise}[Deep research evaluation]
\TODO{}
\end{exercise}

% ──────────────────────────────────────────────────────────────────────────────
\section{Computer-Use \& GUI Agents}
\sectionnote{3 papers.}

\begin{description}
\paper{2026}{OmegaUse}{}{\pBp}
\paper{2026}{GUI-Libra}{}{\pBp}
\paper{2025}{ComputerRL}{}{\pBp}
\end{description}

\begin{exercise}[GUI agent state-action space]
\TODO{Define the state space, action space, and a terminal condition for a GUI agent trained with RL on a web browser.}
\end{exercise}

% ──────────────────────────────────────────────────────────────────────────────
\section{Tool-Use RL}
\sectionnote{3 papers.}

\begin{description}
\paper{2025}{ToolRL}{Reward design textbook; 191 citations}{\pA}
\paper{2025}{Nemotron-Tool-N1}{Binary reward; must-discuss in interviews}{\pA\NV}
\paper{2025}{ReTool}{Cold-start synthesis + outcome RL; 231 citations}{\pAm}
\end{description}

\begin{exercise}[Tool-use reward design]
\TODO{ToolRL uses a structured reward for tool calls.  Write out the reward components; discuss how each discourages reward hacking.}
\end{exercise}

\begin{exercise}[Cold-start synthesis in ReTool]
\TODO{Why does ReTool use a cold-start synthesis phase before RL?  What distribution shift problem does it solve?}
\end{exercise}

\begin{exercise}[Binary vs dense reward]
\TODO{Compare the binary reward in Nemotron-Tool-N1 with a dense shaped reward; give one advantage and one disadvantage of each.}
\end{exercise}

% ──────────────────────────────────────────────────────────────────────────────
\section{Generalist Agent RL Frameworks}
\sectionnote{5 papers.}

\begin{description}
\paper{2026}{RLAnything}{Joint optimisation of env + policy + RM}{\pAm}
\paper{2025}{AGENTRL}{Async + cross-policy sampling}{\pAm}
\paper{2025}{Kimi K2.5 Agent Swarm}{Multi-agent orchestration for agentic tasks}{\pBp}
\paper{2025}{WebAgent-R1}{End-to-end web-agent RL; 73 citations}{\pBp}
\paper{2026}{ProRL Agent}{NVIDIA; rollout-as-a-service for multi-turn agent RL}{\pBp\NV}
\end{description}

\begin{exercise}[Joint optimisation in RLAnything]
\TODO{RLAnything jointly optimises environment, policy, and reward model.  Write the tri-level optimisation objective and identify which level is innermost.}
\end{exercise}

\begin{exercise}[Cross-policy sampling]
\TODO{}
\end{exercise}

\begin{exercise}[Rollout-as-a-service]
\TODO{Describe the rollout-as-a-service abstraction in ProRL Agent; explain how it decouples the policy server from the environment server.}
\end{exercise}

% ──────────────────────────────────────────────────────────────────────────────
\section{Agent RL Challenges \& Self-Evolution}
\sectionnote{4 papers.}

\begin{description}
\paper{2025}{RAGEN}{Echo Trap — why agent RL fails; 148 citations}{\pA}
\paper{2026}{iStar (ICLR 2026)}{Implicit PRM for credit assignment}{\pAm}
\paper{2025}{AgentPRM}{MC rollout actor-critic; 29 citations}{\pBp}
\paper{2025}{Absolute Zero}{Zero-data self-play; 161 citations}{\pAm}
\end{description}

\begin{exercise}[Echo Trap]
\TODO{Define the Echo Trap from RAGEN; explain why it arises specifically in agentic (multi-turn) settings but not in single-turn RL.}
\end{exercise}

\begin{exercise}[Implicit PRM for credit assignment]
\TODO{How does iStar assign per-step credit without explicit step-level annotations?  Sketch the MC estimate used.}
\end{exercise}

\begin{exercise}[Absolute Zero self-play data generation]
\TODO{Describe how Absolute Zero generates its own training data without any human-labelled examples.}
\end{exercise}

\end{document}
